% Établissement du courant. Résolution analytique
Une bobine d'inductance $L$ et de résistance $R$ est placée dans un montage
série comprenant~: un générateur de tension continue $E=\qty{5.0}{\V}$, un
interrupteur $\mathrm{K}$, un conducteur ohmique de résistance $R^{\prime}$. À
la date $t=0$, on ferme l'interrupteur $\mathrm{K}$. On note $r=R+R^{\prime}$.

\begin{enumerate}
    \item Faire un schéma du circuit. Représenter par une flèche le sens
          conventionnel de circulation du courant d'intensité $i$ Représenter
          par des flèches les tensions $u_{L}$ aux bornes de la bobine et
          $u_{R}$ aux bornes du conducteur ohmique, en respectant la convention
          récepteur.
    \item En appliquant la loi d'additivité des tensions, établir une relation
          entre $E$, $u_{R^{\prime}}$ et $u_{L^{\prime}}$.
    \item Établir l'équation différentielle à laquelle obéit $i$. Cette équation
          sera notée~(1).
    \item Une solution de cette équation différentielle est de la forme~:
          $i=c+a\times\mathrm{e}^{bt}$.
          \begin{enumerate}
              \item Déterminer la valeur des constantes $b$ et $c$.
              \item Que vaut l'intensité $i$ du courant à $t=0$~? En déduire la
                    valeur du coefficient $a$.
              \item Exprimer la tension $i$ en fonction de la date $t$.
              \item Vers quelle limite l'intensité tend-elle au bout d'un temps
                    très long~?
          \end{enumerate}
\end{enumerate}


